\twocolumn
[
\begin{@twocolumnfalse}
\maketitle
\begin{abstract}
\vspace*{0.5cm}
\justify
\fontsize{10pt}{10pt}\selectfont

La contaminación atmosférica representa una amenaza creciente para la biodiversidad urbana, pero su efecto sobre la distribución de aves neotropicales ha sido poco estudiado en ciudades latinoamericanas. En este trabajo, analizamos cómo los niveles de material particulado ($\text{PM}_{10}$) y ozono ($\text{O}_3$) se relacionan con la presencia del copetón (\textit{Zonotrichia capensis}) en Bogotá. Para esto, integramos registros de la plataforma de ciencia ciudadana eBird con datos de la Red de Monitoreo de Calidad del Aire de Bogotá (RMCAB). Implementamos un modelo bayesiano de ocupación que nos permite corregir los errores de observación y distinguir entre la ausencia real del ave y el hecho de que simplemente no haya sido vista. Al usar un enfoque probabilístico con distribuciones que se ajustan a los datos, buscamos determinar si la polución reduce significativamente la presencia de esta especie en la capital. Con ello, este trabajo busca responder una pregunta concreta: ¿afecta la contaminación atmosférica urbana la probabilidad de observar al copetón en Bogotá?

\keywordsSpan{Modelo de ocupación, Estadística bayesiana, Pólya-Gamma, Contaminación atmosférica, \textit{Zonotrichia capensis}, eBird, Bogotá}
\end{abstract}
\vspace{0.5cm}



\renewcommand{\abstractname}{ABSTRACT}
\begin{abstract}
\vspace*{0.5cm}
\justify

\textit{Air pollution is a growing threat to urban biodiversity, yet its effect on Neotropical birds remains poorly understood in Latin American cities. In this study, we analyze how particulate matter ($\text{PM}_{10}$) and ozone ($\text{O}_3$) levels affect the occurrence of the Rufous-collared Sparrow (\textit{Zonotrichia capensis}) in Bogotá. We integrated records from the eBird citizen science platform with data from Bogotá's Air Quality Monitoring Network (RMCAB). We implemented a Bayesian occupancy model to separate the actual presence of the bird from observation errors, ensuring that the results account for the uncertainty inherent in citizen science data. By using a data-driven probabilistic approach, we aim to determine if air pollution significantly reduces the presence of this species in the city. This work seeks to answer a concrete question: does urban air pollution affect the probability of observing the Rufous-collared Sparrow in Bogotá?}

\keywordsEng{Occupancy model, Bayesian statistics, Pólya-Gamma, Air pollution, \textit{Zonotrichia capensis}, eBird, Bogotá}
\end{abstract}

\vspace*{0.5cm}  %%% Espacio entre el abstract y las columnas

\end{@twocolumnfalse}
]

